\section{Finite Classification and Completeness}
\label{app:finite-classification}

This appendix gives the finite part of the classification as a mathematical
proof. Its computational input consists of finite integer matrices,
integer vectors, rational inequalities, polynomial factorization, and
rational positive-definiteness certificates; every conclusion below follows
from these data.

\subsection{The finite quotient and its spectral certificates}
\label{app:finite-quotient}

We first make precise why the finite domains used below cover every signing.
In the notation of Section~\ref{sec:switching-reference}, switching along a
Hamilton path gives
\[
 a_0=\cdots=a_{n-2}=1,
 \qquad a_{n-1}=\alpha\in\{\pm1\}.
\]
The switching class is then determined by the triangle-flux word
\(\tau\in\{\pm1\}^n\) and the holonomy \(\alpha\).  On passing to
\[
 Q_i=\tau_i\tau_{i+1},
\]
one obtains the necessary cyclic condition
\begin{equation}\label{eq:appB-Q-product}
 \prod_{i=0}^{n-1}Q_i=1.
\end{equation}
Conversely, every cyclic \(Q\)-word satisfying
\eqref{eq:appB-Q-product}, together with a choice of \(\tau_0\), reconstructs
a unique cyclic lift through \(\tau_{i+1}=Q_i\tau_i\), and either value of
\(\alpha\) then reconstructs a signing.  Thus \((Q,\tau_0,\alpha)\) covers
all switching classes.

For even \(n\), the two choices of \(\tau_0\) need not have the same
unsquared spectrum, but equation~\eqref{eq:lift-negation} gives
\[
 A_{-\tau,\alpha}=-DA_{\tau,\alpha}D,
 \qquad
 A_{-\tau,\alpha}^{\,2}=DA_{\tau,\alpha}^{\,2}D,
 \qquad D=\diag\bigl((-1)^i\bigr)_{i=0}^{n-1}.
\]
They therefore have identical squared spectra and spectral radii. For the
present minimization problem we fix \(\tau_0=1\) after this identification
and retain both holonomies.

Rotation of \(Q\) is induced by relabelling the vertices.  Reflection is
also induced by a graph isomorphism; if the normalization \(\tau_0=1\) is
restored afterward, it may additionally negate the lift.  Hence both
operations preserve the squared spectrum and spectral radius. Passing to a
dihedral representative retains every possible value of \(\rho(A)\), with
the two values of \(\alpha\) kept separate.

We use two elementary certificate principles. Since \(A\) is real
symmetric,
\begin{equation}\label{eq:appB-rayleigh}
 \rho(A)^2=\lambda_{\max}(A^2)
 =\max_{v\ne0}\frac{\lVert Av\rVert^2}{\lVert v\rVert^2}.
\end{equation}
Thus an integer vector \(v\ne0\) for which
\begin{equation}\label{eq:appB-rayleigh-lower}
 \lVert Av\rVert^2\geq\theta_n\lVert v\rVert^2
\end{equation}
proves \(\rho(A)^2\geq\theta_n\).  The comparison in
\eqref{eq:appB-rayleigh-lower} is certified: \(\theta_n\) is selected as an
isolated real root of its integer minimal polynomial, so the sign of the
algebraic difference is decided by rational interval isolation.

In the opposite direction, if positive integers \(p,q\) satisfy
\begin{equation}\label{eq:appB-pd-certificate}
 pI-qA^2\succ0,
\end{equation}
then every eigenvalue \(\lambda\) of \(A\) satisfies
\(p-q\lambda^2>0\), and hence
\begin{equation}\label{eq:appB-pd-consequence}
 \rho(A)^2<\frac pq.
\end{equation}
Positive definiteness in \eqref{eq:appB-pd-certificate} is decided over
\(\mathbb Q\), either from strictly positive fraction-free leading principal
minors or from a rational factorization \(LDL^{\mathsf T}\) with every diagonal entry of
\(D\) positive.

Finally, for every even \(n\),
\begin{equation}\label{eq:appB-cosine-bound}
 \theta_n
 =4+2\cos\frac{2\pi}{n}+2\cos\frac{4\pi}{n}
 >8-\frac{20\pi^2}{n^2}
 >8-\frac{200}{n^2}.
\end{equation}
The first strict inequality follows from
\(\cos x>1-x^2/2\) for nonzero \(x\), and the second from \(\pi^2<10\).

\subsection{Complete decisions through order thirty}
\label{app:orders-8-30}

For \(8\leq n\leq20\), direct enumeration uses one representative of each
switching class.  Since \(G_n\) is connected and has \(2n\) edges, the
number of switching classes is
\[
 \frac{2^{2n}}{2^{n-1}}=2^{n+1}.
\]
For \(n=22\), the same space is represented in the \((Q,\alpha)\) quotient
described above.  At \(n=24,26,28,30\), cyclic \(Q\)-words are generated by
weight and then reduced under rotations and reflections.  Summing the
represented dihedral-orbit sizes and retaining both values of \(\alpha\)
gives \(2^n\) reduced \((Q,\alpha)\) states.  Each such state represents the
two choices \(\tau_0=\pm1\), so the weighted switching-class total is
\(2^{n+1}\).  Their spectral decision may be shared only because the two
lifts have identical squared spectra.  Thus these are reorganizations of the
full finite domain, not restricted searches.

Every represented class is discharged in one of two ways.  A nonthreshold
class carries a primitive integer vector satisfying
\eqref{eq:appB-rayleigh-lower}. In a threshold class, the
characteristic polynomial of \(A^2\) contains the minimal polynomial of
\(\theta_n\), and isolation of all real roots proves that this factor
supplies the largest squared eigenvalue.  Hence such a class has
\(\rho(A)^2=\theta_n\).  The complete finite decisions are summarized in
Table~\ref{tab:appB-small-orders}.

\begin{table}[t]
\centering
\caption{Exhaustive decisions for even orders at most thirty. In the
last five rows, a quotient record can represent an entire dihedral orbit.}
\label{tab:appB-small-orders}
\begin{tabular}{@{}r r r l@{}}
\toprule
\(n\)&Switching classes&Nonthreshold records&Threshold closure\\
\midrule
8  &512          &510       &2 threshold equalities\\
10 &2,048        &2,046     &2 threshold equalities\\
12 &8,192        &8,190     &2 threshold equalities\\
14 &32,768       &32,766    &2 threshold equalities\\
16 &131,072      &131,070   &2 threshold equalities\\
18 &524,288      &524,286   &2 threshold equalities\\
20 &2,097,152    &2,097,150 &2 threshold equalities\\
22 &8,388,608    &97,467    &threshold quotient\\
24 &33,554,432   &353,811   &threshold quotient\\
26 &134,217,728  &1,299,063 &threshold quotient\\
28 &536,870,912  &4,810,471 &threshold quotient\\
30 &2,147,483,648&17,929,599&threshold quotient\\
\bottomrule
\end{tabular}
\end{table}

It follows from the exhaustive alternative above that every signing at these
orders satisfies \(\rho(A)^2\geq\theta_n\).  Proposition
\ref{prop:candidate-attainment} supplies, at every even order, a signing with
\(\rho(A)^2=\theta_n\).  Therefore
\begin{equation}\label{eq:appB-small-equality}
 m_n=\rho_-(n)
 \qquad(8\leq n\leq30,\ n\ \text{even}).
\end{equation}

\subsection{The counterexamples at orders thirty-two and forty}
\label{app:orders-32-40}

At order \(32\), take \(\alpha=+1\) and the triangle-flux word
\[
 \tau=(+,+,-,+,-,-,+,-)^4.
\]
The Hamilton-gauge reconstruction gives an integer symmetric signed
adjacency matrix \(A_{32}\). Fraction-free elimination, followed by
an independent rational \(LDL^{\mathsf T}\) factorization of the same
reconstructed matrix, gives
\begin{equation}\label{eq:appB-n32-pd}
 1561I_{32}-200A_{32}^2\succ0.
\end{equation}
By \eqref{eq:appB-pd-consequence},
\(\rho(A_{32})^2<1561/200\). Root isolation for \(\theta_{32}\)
gives
\[
 \frac{1561}{200}
 <\frac{11896117236720419}{1523321182060814}
 <\theta_{32}.
\]
Consequently \(m_{32}<\rho_-(32)\).  The matrix inequality
\eqref{eq:appB-n32-pd} is an upper certificate for this one explicit
signing; it neither computes the value of \(m_{32}\) nor asserts that
the displayed signing is the unique minimizer.

At order \(40\), encode \(Q_i=+1\) by the digit \(1\) and take
\[
 Q=1000100010001000100010001000100010001000,
 \qquad \alpha=-1,
\]
with \(\tau_0=1\).  The recurrence \(\tau_{i+1}=Q_i\tau_i\) and Hamilton
gauge reconstruct a unique integer symmetric matrix \(A_{40}\).
Rational elimination gives forty positive pivots in
\begin{equation}\label{eq:appB-n40-pd}
 15541I_{40}-2000A_{40}^2,
\end{equation}
so
\[
 \rho(A_{40})^2<\frac{15541}{2000}.
\]
Moreover,
\[
 \frac{15541}{2000}<\frac{63}{8}
 =8-\frac{200}{40^2}<\theta_{40},
\]
where the first difference is \(209/2000\) and the last inequality is
\eqref{eq:appB-cosine-bound}.  Hence \(m_{40}<\rho_-(40)\).  This
full-matrix upper certificate is logically independent of the lower-bound
classification at the six intervening orders.

\subsection{Local exclusion at the six recovery orders}
\label{app:local-exclusion}

We now prove the finite-state reduction used for
\[
 \mathcal N=\{34,36,38,42,44,46\}.
\]
For each \(n\in\mathcal N\), cyclotomic elimination gives the minimal
polynomial of \(\theta_n\).  The value \(\theta_n\) is its rightmost real
root: if \(c\) is a conjugate of \(\cos(2\pi/n)\), then the corresponding
conjugate is
\[
 f(c)=2+2c+4c^2.
\]
The function \(f\) is increasing on \([-1/4,1]\), while \(f(c)\leq4\) on
\([-1,-1/4]\); the smallest nonzero angular distance among the cyclotomic
conjugates is \(2\pi/n\). Sturm counts therefore select strict
rational endpoints
\begin{equation}\label{eq:appB-isolating-interval}
 a_n<\theta_n<b_n.
\end{equation}

Let \(S=\{0,\ldots,L-1\}\), where
\begin{equation}\label{eq:appB-support-lengths}
 L_{34}=12,\qquad L_{36}=13,\qquad
 L_{38}=L_{42}=L_{44}=L_{46}=14.
\end{equation}
Because these supports and their range-two collars are proper arcs of the
corresponding cycles, the holonomy seam can be moved outside the collar by
switching.  For a vector supported on \(S\), equation~\eqref{eq:A-tau}
defines an integer map
\[
 C_W:\mathbb R^L\longrightarrow\mathbb R^{L+4},
\]
where the length-\((L+1)\) word
\[
 W=(Q_{-2},Q_{-1},\ldots,Q_{L-2})
\]
determines all required \(\tau\)-coefficients up to simultaneous negation.
That negation does not change the squared spectrum.  If \(P_S\) denotes
compression to \(S\), then
\begin{equation}\label{eq:appB-local-compression}
 M_W:=C_W^{\mathsf T}C_W=P_SA^2P_S.
\end{equation}

\begin{lemma}[Local exclusion]\label{lem:appB-local-exclusion}
If an integer vector \(v\ne0\) satisfies
\[
 v^{\mathsf T}M_Wv>b_n v^{\mathsf T}v,
\]
then every cyclic signing containing \(W\) satisfies
\(\rho(A)^2>\theta_n\).
\end{lemma}

\begin{proof}
Extend \(v\) by zero off \(S\).  Equations
\eqref{eq:appB-rayleigh} and \eqref{eq:appB-local-compression} give
\[
 \rho(A)^2\geq
 \frac{v^{\mathsf T}M_Wv}{v^{\mathsf T}v}
 >b_n>\theta_n.
\]
The argument is unaffected by the location of the window or the holonomy
seam.
\end{proof}

Every one of the \(2^{L+1}\) windows is decided in certified arithmetic. A
surviving window satisfies
\[
 a_nI-M_W\succ0,
\]
as proved by the fraction-free Sylvester criterion; an excluded window carries an
integer vector satisfying Lemma~\ref{lem:appB-local-exclusion}.  The two
lists form a partition of all windows.  Thus a signing is either already
excluded by a local Rayleigh certificate or every one of its cyclic windows
belongs to the finite surviving set.

\subsection{Parity-lifted overlap closure}
\label{app:parity-overlap}

Encode \(Q_i=+1\) by \(b_i=1\) and \(Q_i=-1\) by \(b_i=0\).  Since \(n\)
is even, condition~\eqref{eq:appB-Q-product} becomes
\begin{equation}\label{eq:appB-bit-parity}
 \sum_{i=0}^{n-1}b_i\equiv0\pmod2.
\end{equation}
Let \(\mathcal W_n\) be the surviving binary words of length \(L+1\).
Construct a directed overlap graph with length-\(L\) states: a word
\((b_0,\ldots,b_L)\in\mathcal W_n\) is the edge
\[
 (b_0,\ldots,b_{L-1})
 \longrightarrow(b_1,\ldots,b_L).
\]
Lift every state to \((s,\varepsilon)\), with
\(\varepsilon\in\mathbb Z/2\mathbb Z\), and lift the edge by
\begin{equation}\label{eq:appB-parity-transition}
 (s,\varepsilon)\longrightarrow
 (s',\varepsilon+b_L).
\end{equation}

\begin{lemma}[Soundness and completeness]
\label{lem:appB-debruijn-completeness}
Length-\(n\) closed walks in the parity-lifted overlap graph are in
bijection with rooted cyclic \(Q\)-words satisfying
\eqref{eq:appB-Q-product} and containing no excluded window.
\end{lemma}

\begin{proof}
Along a walk, consecutive edges overlap in exactly \(L\) symbols, so their
appended labels form a word of length \(n\).  Closure of the base state
forces the last \(L\) overlaps with the first \(L\); hence every wraparound
window also survives.  Closure of the parity coordinate in
\eqref{eq:appB-parity-transition} says that the sum of all appended bits is
even, which is precisely \eqref{eq:appB-bit-parity}.  This constructs a
legal cyclic \(Q\)-word and proves soundness.

Conversely, start at any root of a legal cyclic \(Q\)-word having no
excluded window and read its symbols consecutively.  Each length-\((L+1)\)
window is an edge of the overlap graph, successive windows have the required
overlap, cyclicity closes the base state, and
\eqref{eq:appB-bit-parity} closes the parity state.  This produces a
length-\(n\) closed walk and proves completeness.
\end{proof}

After enumerating these closed walks, rotation and reflection identify
different roots and orientations of the same cyclic configuration.  By the
spectral equivalences proved in Subsection~\ref{app:finite-quotient}, taking
one dihedral representative is exhaustive.  For every resulting \(Q\)-word,
both \(\alpha=-1\) and \(\alpha=+1\) are retained.  The two \(\tau\)-lifts
are already covered by their identical squared spectra.  Consequently the
terminal \((Q,\alpha)\)-records cover every signing that was not locally
excluded.

The finite counts are given in Table~\ref{tab:appB-recovery-orders}.
The rooted even and odd columns count closed walks by total bit parity before
the cyclic and dihedral reductions; only the even walks are globally
liftable.  Different roots can represent the same cyclic word.

\begin{table}[t]
\centering
\caption{Finite-state closure at the six recovery orders.}
\label{tab:appB-recovery-orders}
\small
\setlength{\tabcolsep}{4pt}
\begin{tabular}{@{}r r r r r r r r@{}}
\toprule
\(n\)&\(L\)&Surviving&States&Rooted even&Rooted odd&Canonical \(Q\)&Terminals\\
\midrule
34&12&124&92 &1  &0  &1 &2\\
36&13&128&92 &1  &4  &1 &2\\
38&14&184&132&77 &38 &3 &6\\
42&14&232&166&337&392&7 &14\\
44&14&240&171&353&620&10&20\\
46&14&240&171&599&690&10&20\\
\bottomrule
\end{tabular}
\end{table}

There are therefore
\begin{equation}\label{eq:appB-terminal-total}
 2(1+1+3+7+10+10)=64
\end{equation}
terminal \((Q,\alpha)\)-records.  Six of them have the all-negative
\(Q\)-word and \(\alpha=-1\).  For each of these, the minimal polynomial of
\(\theta_n\) divides \(\det(xI-A^2)\), and root isolation shows that
this factor gives the largest squared eigenvalue.  Thus these six terminals
satisfy \(\rho(A)^2=\theta_n\).  Each of the other \(58\) terminals carries
a primitive integer vector \(v\) with
\[
 \frac{\lVert Av\rVert^2}{\lVert v\rVert^2}>b_n>\theta_n.
\]
All \(64\) terminals are therefore decided.

Every signing at one of the six recovery orders either contains an excluded
window, in which case Lemma~\ref{lem:appB-local-exclusion} applies, or gives
a parity-even closed walk, in which case
Lemma~\ref{lem:appB-debruijn-completeness} places it in one of the resolved
terminals.  Hence
\[
 \rho(A)^2\geq\theta_n
 \qquad(n=34,36,38,42,44,46).
\]
Combining this universal lower bound with Proposition
\ref{prop:candidate-attainment} yields
\begin{equation}\label{eq:appB-recovery-equality}
 m_n=\rho_-(n)
 \qquad(n=34,36,38,42,44,46).
\end{equation}

\subsection{The finite bridge from forty-eight to two hundred thirty-eight}
\label{app:finite-bridge}

It remains to record the finite boundary between the isolated
small orders and the analytic tail.  The even interval
\[
 48\leq n<240
\]
contains \((238-48)/2+1=96\) orders. At each such order, define
\(A_n\) by the deterministic residue family
\[
\begin{array}{c|l|c}
n\bmod8&\text{flux construction}&\alpha\\
\hline
0&\text{period-eight repetition}&-1\\
2&\text{one G6 phase slip}&+1\\
4&\text{two balanced G6 phase slips}&-1\\
6&\text{three balanced G6 phase slips}&+1.
\end{array}
\]
The gap words are those of Section~\ref{sec:finite-rings}; the cyclic lift
condition is checked before Hamilton gauge, so every row determines a signed
cycle square.

For every even \(n\) in this interval, the finite certificate supplies
positive integers \(p_n,q_n\) such that
\begin{equation}\label{eq:appB-bridge-certificate}
 p_nI-q_nA_n^2\succ0,
 \qquad
 \frac{p_n}{q_n}<8-\frac{200}{n^2}.
\end{equation}
Each matrix \(A_n\) is reconstructed from its displayed residue rule, and
each first inequality is verified by a rational
\(LDL^{\mathsf T}\) factorization with positive diagonal.  Thus the finite
statement covers all \(96\) consecutive even orders. Equations
\eqref{eq:appB-pd-consequence}, \eqref{eq:appB-bridge-certificate}, and
\eqref{eq:appB-cosine-bound} give, separately for every row,
\[
 \rho(A_n)^2
 <\frac{p_n}{q_n}
 <8-\frac{200}{n^2}
 <\theta_n.
\]
Therefore
\begin{equation}\label{eq:appB-bridge-failure}
 m_n<\rho_-(n)
 \qquad(48\leq n<240,\ n\ \text{even}).
\end{equation}
The endpoint \(238\) is included in this finite bridge, and the residue-wise
IMS propagation in Section~\ref{sec:finite-rings} begins at the next even
order, \(240\).

Equations \eqref{eq:appB-small-equality},
\eqref{eq:appB-recovery-equality}, the strict witnesses at \(32\) and
\(40\), and \eqref{eq:appB-bridge-failure} are the finite assertions
used in the proof of Theorem~\ref{thm:complete-classification}.
